
%%%%

%%%   Самарский государственный технический университет

%%%   Кафедра прикладной математики и информатики

%%%

%%%   Преамбула для статьи в журнал "Вестник Самарского государственного

%%%   технического университета. Серия: «Физико-математические науки»"

%%%   Ориентирована под MikTEx 2.4 for Win32

%%%

%%%   Хотя сам журнал собирается под GNU/Linux на дистрибутиве TeXLive



\documentclass[11pt,a4paper,twoside]{article}



\usepackage[cp1251]{inputenc}

\usepackage[T2A]{fontenc}

\usepackage[english,russian]{babel}

\usepackage{samgtubib}

\usepackage[dvips]{graphicx}

\usepackage[multiple]{footmisc}



\usepackage{samgtu}

\renewcommand{\VestnikYear}{2009} % Год издания Вестника

\renewcommand{\VestnikNumber}{2\,(19)} % Номер Вестника

\renewcommand{\VestnikMonth}{Ноябрь} % Месяц выхода Вестника

\renewcommand{\VestnikData}{01.11.09} % Дата подписания Вестника в печать

\renewcommand{\VestnikUPL}{26,64} % Усл. печ. л.

\renewcommand{\VestnikUIL}{25,73} % Уч.-изд. л.

\renewcommand{\VestnikRegNumber}{479/09} % Регистрационный номер

\renewcommand{\VestnikOrder}{994} % Номер заказа








%\begin{document}
\renewcommand{\firstpage}{180}
\renewcommand{\lastpage}{183}



\udk{[544.344.9:544.015.33](043.3)}

\title{Использование программного комплекса <<Dif Pro Generator>> 
для~исследования и~уточнения характеристик эвтектик тр\"eхкомпонентных взаимных 
систем}{Использование программного комплекса <<Dif Pro Generator>> \dots} % 
\engtitle{Application of the programm complex ``Dif Pro Generator'' for of~three--component mutual systems eutectic characteristic research and spesification} 

\author{А.\,С.~Трунин, О.\,Е.~Моргунова, С.\,В.~Горбачёв}
       {Т\,р\,у\,н\,и\,н~А.\,С.,  М\,о\,р\,г\,у\,н\,о\,в\,а~О.\,Е.,  Г\,о\,р\,б\,а\,ч\,ё\,в~С.\,В.}% Автор(ы) для Оглавления и колонтитула         
\engauthor{A.\,S.~Trunin, O.\,Y.~Morgunova, S.\,V.~Gorbachev}



\date{02.10.2006} % Дата поступления статьи в редакцию.


\thanks{}

\address{\samgtu}
\engaddress{Samara State Technical University, Samara, Russia}

\email{olvale@mail.ru}

\workabstract{Рассмотрена возможность применения разработанного программного комплекса <<Dif Pro Generator>> для моделирования характеристик эвтектик в качестве самостоятельного метода исследования трёхкомпонентных взаимных физико"=химических систем, а~также экспресс"=метода для выбора наиболее достоверной информации при наличии противоречивых данных о~составах и~температурах нонвариантных точек у различных исследователей.
}

\engabstract{In the article the authors consider possibility of worked out program complex ``Dif Pro Generator'' application for eutectic characteristics modeling as independent method of three-component mutual systems research and as express--method for more reliable information selection, if data of different investigators about non-variant points compositions and temperature are contradict.}

\maketitle

Компьютерное моделирование является наиболее перспективным направлением 
развития методологии исследования многокомпонентных систем на современном 
этапе. В~работе [1] предложен метод компьютерного моделирования и~расчёта 
эвтектических характеристик трёхкомпонентных систем с~помощью разработанного 
и~реализованного программного комплекса $\tt Dif\,\,Pro\,\,Generator$~[2]. Он 
позволяет получать состав и температуру эвтектики трёхкомпонентной системы 
с~точностью до 5\,{\%}, соизмеримой с~экспериментальной, без проведения 
эксперимента. Однако в~работе [1] рассматривается лишь моделирование чисто 
эвтектических трёхкомпонентных систем. Развитие идеи предполагает 
распространение её на системы с~разнообразным морфологическим типом 
взаимодействия компонентов (обменные реакции, комплексообразования 
конгруэнтного и~инконгруэнтного характера, твёрдых растворов различной 
степени устойчивости и~др.). 

Предлагается рассмотреть возможность применения программного комплекса для 
получения состава и~температуры эвтектик трёхкомпонентных 
необратимо"=взаимных физико"=химических систем.

Особенностью данного вида систем является наличие обменных реакций, поэтому 
при их исследовании возникает дополнительная задача разбиения системы на 
фазовые единичные блоки (ФЕБы)~[3]. Процедура дифференциации предполагает 
наличие дополнительных сведений (энтальпия образования $\Delta  H$) об 
однокомпонентных элементах огранения.

Рассмотрим процедуру реализации предлагаемого подхода на примере системы Na,\,K\,/\!/\,F,\,Cl (см. рис.~1). 
Она была изучена ранее различными авторами~[4]. Из 
табл.~1 видно, что характеристики эвтектик данной тройной взаимной системы 
значительно отличаются друг от друга, как по температуре, так и~по составам.



\begin{table}[b!]
\renewcommand{\bfdefault}{b}
\caption{\small \bfseries Экспериментальные данные по характеристикам эвтектик системы Na,\,K\,/\!/\,F,\,Cl~[4]}
\vspace{-3mm} 
\noindent
\begin{tabular}{l|l|l|l|l|l|l|l}
\hline
\multicolumn{1}{c|}{Обозна-} & \multicolumn{1}{c|}{Эвтек-} & \multicolumn{4}{c|}{Содержание компо-} &  &  \\ 
\multicolumn{1}{c|}{чение на } & \multicolumn{1}{c|}{тика,} & \multicolumn{4}{c|}{нентов, экв. \%} & \multicolumn{1}{c|}{Метод исследования} & \multicolumn{1}{c}{\totop{Авторы }} \\ 
\cline{3-6}
\multicolumn{1}{c|}{рис.~1} & \multicolumn{1}{c|}{$t$,  $\rm^\circ C$} & \multicolumn{1}{c|}{KCl} & \multicolumn{1}{c|}{KF} & \multicolumn{1}{c|}{NaCl} & \multicolumn{1}{c|}{NaF} &  & \multicolumn{1}{c}{\totop{исследования}} \\ 
\hline
\multicolumn{1}{c|}{$E_1$} & \multicolumn{1}{c|}{570} & \multicolumn{1}{c|}{46,5} & \multicolumn{1}{c|}{39,5} & \multicolumn{1}{c|}{---} & \multicolumn{1}{c|}{14} & Визуально"=политер- & \multicolumn{1}{c}{} \\ 
\multicolumn{1}{c|}{$E_2$} & \multicolumn{1}{c|}{606} & \multicolumn{1}{c|}{22,5} & \multicolumn{1}{c|}{---} & \multicolumn{1}{c|}{60,5} & \multicolumn{1}{c|}{17} & мический (ВПМ) & \multicolumn{1}{c}{Поляков В.\,Д. } \\ 
\hline
\multicolumn{1}{c|}{$E_1$} & \multicolumn{1}{c|}{---} & \multicolumn{1}{c|}{---} & \multicolumn{1}{c|}{---} & \multicolumn{1}{c|}{---} & \multicolumn{1}{c|}{---} & ВПМ, кривые  охлаж- & \multicolumn{1}{c}{Алабышев А.\,Ф., } \\ 
\multicolumn{1}{c|}{$E_2$} & \multicolumn{1}{c|}{615--603} & \multicolumn{1}{c|}{42,5} & \multicolumn{1}{c|}{---} & \multicolumn{1}{c|}{42,5} & \multicolumn{1}{c|}{15} & дения & \multicolumn{1}{c}{Купперберг Л.С. } \\ 
\hline
\multicolumn{1}{c|}{$E_1$} & \multicolumn{1}{c|}{582} & \multicolumn{1}{c|}{---} & \multicolumn{1}{c|}{---} & \multicolumn{1}{c|}{---} & \multicolumn{1}{c|}{---} &  & \multicolumn{1}{c}{Chretien A., } \\ 
\multicolumn{1}{c|}{$E_2$} & \multicolumn{1}{c|}{612} & \multicolumn{1}{c|}{---} & \multicolumn{1}{c|}{---} & \multicolumn{1}{c|}{---} & \multicolumn{1}{c|}{---} & \totop{Кривые охлаждения} & \multicolumn{1}{c}{Silber P., Ichaque M. \!\!\!\!\!\!\!} \\ 
\hline
\multicolumn{1}{c|}{} & \multicolumn{1}{c|}{} & \multicolumn{1}{c|}{} & \multicolumn{1}{c|}{} & \multicolumn{1}{c|}{} & \multicolumn{1}{c|}{} & Кривые охлаждения,  & \multicolumn{1}{c}{} \\ 
\multicolumn{1}{c|}{\totop{$E_1$}} & \multicolumn{1}{c|}{\totop{---}} & \multicolumn{1}{c|}{\totop{---}} & \multicolumn{1}{c|}{\totop{---}} & \multicolumn{1}{c|}{\totop{---}} & \multicolumn{1}{c|}{\totop{---}} & микроскопическое  иссле- & \multicolumn{1}{c}{Рассонская И.\,С., } \\ 
\multicolumn{1}{c|}{} & \multicolumn{1}{c|}{} & \multicolumn{1}{c|}{} & \multicolumn{1}{c|}{} & \multicolumn{1}{c|}{} & \multicolumn{1}{c|}{} & дование твёрдых шлифов  & \multicolumn{1}{c}{Бергман А.\,Г. } \\ 
\multicolumn{1}{c|}{\totop{$E_2$}} & \multicolumn{1}{c|}{\totop{590}} & \multicolumn{1}{c|}{\totop{---}} & \multicolumn{1}{c|}{\totop{---}} & \multicolumn{1}{c|}{\totop{---}} & \multicolumn{1}{c|}{\totop{---}} & отожжённых сплавов & \multicolumn{1}{c}{} \\ 
\hline
\end{tabular}
\end{table}

Использование идеологии моделирования [1] с~применением программного 
комплекса $\tt Dif\,\,Pro\,\,Ge$- $\tt nerator$~[2] для исследования эвтектик 
трёхкомпонентных взаимных систем и~в~качестве критерия выбора из нескольких 
противоречивых экспериментальных данных различных исследователей 
предполагает осуществление ряда этапов.


\begin{wrapfigure}{O}{.46\textwidth}
\centering
\begin{minipage}{.45\textwidth}
\centerline{\includegraphics[width=.97\textwidth]{./00Pic/trunin1.eps}}
\vspace{.1cm}\centering
\parbox[b]{.8\textwidth}{\caption{Диаграмма плавкости системы Na,\,K\,/\!/ F,\,Cl~[4]}}
\end{minipage}
\end{wrapfigure}

1. Анализ топологии системы с целью выявления возможности применения ЭГ. Как видно из рис.~1, стабильная диагональ NaF--KCl делит систему на два вторичных треугольника "--- симплекса: NaF--KCl--KF и NaF--NaCl--KCl. Последний содержит непрерывный ряд твёрдых растворов в~двухкомпонентной системе из хлоридов натрия и~калия, поэтому предлагается провести апробацию работы ЭГ на симплексе NaF--KCl--KF. 

2. Анализ метрики системы с целью выявления возможности применения ЭГ. 
Информация об элементах огранения (одно- и~двухкомпонентных системах) взята 
из [5,\,6]. 

Температуры плавления однокомпонентных систем "--- элементов огранения: 
NaF "--- 995\,$\rm^\circ C$, \linebreak KCl "--- 775\,$\rm^\circ C$, KF "--- 858\,$\rm^\circ C$~[5].

Эвтектика системы NaF--KCl "---  75\,{\%} экв. KCl при 650\,$\rm^\circ C$~[6],
KCl--KF "---  55\,{\%} экв. KCl при 606\,$\rm^\circ C$~[6],
NaF--KF "---  40\,{\%} экв. NaF при 710\,$\rm^\circ C$~[6].

Из приведённых данных о двухкомпонентных элементах огранения и~стабильной 
диагонали системы Na,\,K\,/\!/\,F,\,Cl видно, что все они являются эвтектическими. 
Таким образом, с~достаточной степенью достоверности можно утверждать, что 
вторичный треугольник "--- система NaF--KCl--KF является эвтектической. 
Следовательно, для определения температуры и состава эвтектики симплекса 
тройной взаимной системы возможно использование компьютерного моделирования~[1]. 

3. Расчёт состава и температуры тройной эвтектики системы осуществляется с~использованием программного комплекса $\tt Dif\,\,Pro\,\,Generator$ [2]. Данные компьютерного моделирования приведены в~табл.~2.

\begin{table}[b!]
\vspace{-5mm}
\centering
\begin{minipage}{13cm}
\caption{\small \bfseries Расчётные характеристики эвтектики стабильного треугольника
NaF--KCl--KF трёхкомпонентной взаимной системы Na,\,K\,//\,F,\,Cl}
\vspace{-3mm} 
\noindent
\begin{tabular}{l|l|l|lll}
\hline
\multicolumn{1}{c|}{Обозначение } & \multicolumn{1}{c|}{Температура} & \multicolumn{4}{c}{Содержание компонентов, экв. \%} \\ 
\cline{3-6}
\multicolumn{1}{c|}{на рис. 1} & \multicolumn{1}{c|}{эвтетики, $t$\,$\rm^\circ C$} & \multicolumn{1}{c|}{~~~KCl~~~} & \multicolumn{1}{c|}{~~~KF~~~} & \multicolumn{1}{c|}{~~~NaCl~~~} & \multicolumn{1}{c}{~~~NaF~~~} \\ 
\hline
\multicolumn{1}{c|}{$E_1$} & \multicolumn{1}{c|}{575{,}1} & \multicolumn{1}{c|}{46{,}8} & \multicolumn{1}{c|}{36{,}8} & \multicolumn{1}{c|}{---} & \multicolumn{1}{c}{16{,}4} \\ 
\hline
 \end{tabular}
 \end{minipage}
\renewcommand{\bfdefault}{bx}
 \end{table}



4. Сравнение результатов моделирования характеристик эвтектик на электронном генераторе фазовых диаграмм и экспериментальных данных.

Анализ данных, приведённых в табл. 1 и~2 показал, что в фазовом 
треугольнике NaF--KCl--KF расчётная температура эвтектики $E_1$ 
находится в~пределах экспериментальных значений авторов Полякова~В.\,Д. и~Chretien\,A., Silber\,P., Ichaque\,M., но ближе к~результатам Полякова~В.\,Д. 
Среднеквадратичное значение погрешности расчёта состава эвтектики $E_{1}$ 
относительно эксперимента последнего равно 1{,}8\,{\%}. Такая величина 
погрешности находится в~пределах допустимой экспериментальной (5\,{\%}) [1]. 


Для оценки полученных результатов авторами дополнительно был проведён 
единичный подтверждающий эксперимент на установке дифференциального 
термического анализа.


На рис. 2 и~3 представлены кривые охлаждения эвтектики системы 
KF--KCl--NaF. Состав на рис.~2 рассчитан на ЭГ, а~на рис.~3 определён 
экспериментально исследователем Поляковым~В.\,Д.~[4]. На термограммах виден 
пик кристаллизации эвтектики с~температурой начала кристаллизации эвтектики 
$t_{E_1}=590{,}6$\,$\rm^\circ C$, в~отличие от рис.~2, где имеется фазовый переход при 
$t_{x}=609{,}7$\,$\rm^\circ C$. В~данном случае по величине практически 
незначительного второго пика, с~учётом того, что он не индивидуален, можно 
констатировать о~незначительной величие отклонения от экспериментально 
определённого состава. Это также позволяет сделать заключение о~близости 
расчётного и~экспериментального значения характеристик тройной эвтектики, 
кристаллизующейся в~подсистеме KF--KCl--NaF. Расхождение экспериментально 
полученной температуры эвтектики от данных~[4] и~результатов расчёта на ЭГ 
находится в~пределах 3\,{\%}, что является допустимым.


%\begin{figure}[t!]
\begin{wrapfigure}{O}{.7\textwidth}
\centering
\begin{minipage}{.7\textwidth}
\centerline{\includegraphics[width=.97\textwidth]{./00Pic/trunin2-3.eps}}

\hfill
\begin{minipage} [b]{0.45\textwidth}
\caption{Термическая кривая охлаждения расчётного состава эвтектики стабильной подсистемы KF--KCl--NaF}
\end{minipage}
\hspace{5mm} 
\begin{minipage} [b]{0.45\textwidth}
\caption{Термическая кривая охлаждения эвтектического состава~[5] стабильной подсистемы KF--KCl--NaF }
\end{minipage}
\hfil

\end{minipage}
\end{wrapfigure}


Таким образом, результаты компьютерного моделирования эвтектических 
характеристик стабильного треугольника свидетельствуют о~возможности 
применения электронного генератора фазовых диаграмм в~качестве 
самостоятельного метода для исследования эвтектических трёхкомпонентных 
не\-об\-ратимо"=взаимных систем. Это позволяет сократить количество необходимого 
эксперимента до единичного подтверждающего, а~в~некоторых случаях и~вовсе 
отказаться от проведения эксперимента, т.\,к. погрешность моделирования 
находится в~пределах допустимой экспериментальной. 

Применение ЭГ возможно и в качестве экспертного метода при наличии 
противоречивых данных о~составах и~температурах нонвариантных точек 
у~различных исследователей. Достаточно часто в справочной литературе 
встречаются расхождения у~различных авторов по экспериментальным данным 
одной и~той же трёхкомпонентной эвтектики. Поэтому в~практике 
физико"=химических исследований возникает необходимость выбирать из 
нескольких данных наиболее верные. Методы математической статистики [7,~8] 
дают усреднённые значения, которые могут значительно отличаться от 
истинных.

При необходимости получения более точных данных проводятся тестовые 
экспериментальные проверки всех указанных в~справочниках эвтектик 
с~использованием визуально"=политермического [9] или дифференциального 
термического анализа~[10]. При этом собственный эксперимент принимается за 
основу и~используется в~дальнейших исследованиях. Очевидно, что проведение 
дополнительного эксперимента требует затрат времени и~средств, необходимых 
для подобного тестирования, и,~в~целом, удорожает исследования. 

Предлагается использовать ЭГ в~качестве экспертной системы и~выбирать из 
ряда противоречивых данных те, которые наиболее близки к~расчётным. 
Единственным требованием для моделирования является высокая точность входных 
данных или, как принято их называть, элементов огранения низшей мерности.

\smallskip
\noindent{\bf Выводы}


1. Погрешность моделирования состава и температуры эвтектики вторичного 
треугольника KF--KCl--NaF системы Na,\,K\,/\!/\,F,\,Cl находится в~пределах 3\,{\%}. 
Это значение находится в~пределах принятой допустимой экспериментальной 
погрешности (5\,{\%}), что говорит о~высокой точности моделирования в~данном 
случае и~позволяет предположить возможность практического применения метода~[1] 
не только для изучения эвтектических трёхкомпонентных систем, но и~для 
более сложных систем с~наличием реакций обмена.

2. Использование автоматизированного комплекса $\tt Dif\,\,Pro\,\,Generator$ в~качестве 
экспресс"=метода определения эвтектических характеристик трёхкомпонентных 
систем для выявления наиболее достоверных данных из нескольких вариантов 
характеристик тройных эвтектик в~системах, исследованных в~разное время, 
разными методами и~разными авторами, позволит избежать ошибок и~неточностей 
при исследовании диаграмм состояния трёх- и~более компонентных систем 
экспериментальными и~расчётными методами.

3. Проведённые собственные эксперименты доказали, что достоверными результатами 
исследований эвтектических характеристик являются те, где данные 
моделирования и~эксперимента совпадают или близки.



\begin{thebibliography}{10}

\item \textit{Моргунова, О.\,Е.} Электронный генератор фазовых диаграмм физико"=химических систем [Текст]~/ О.\,Е.~Моргунова, А.\,С.~Трунин.
"--- Самара: СамГТУ, 2005. "--- 133~с.


\item \textit{Чуваков, А.\,В.} 
Программный комплекс <<Dif Pro Generator>> (автоматизированный программный комплекс исследования четырёхкомпонентных взаимных систем)~/ 
А.\,В.~Чуваков, В.\,А.~Лукиных, Н.\,В.~Котляров и др. %Трунин А.С., Климова М.В., Моргунова О.Е., Будкин А.В.
"--- Зарегистрировано в~ОФАП 28.09.2005, № 5180. Код программы по ЕСПД 02068396.00008--01. 


\item \textit{Трунин, А.\,С.} Комплексная методология исследования многокомпонентных систем [Текст]~/ А.\,С.~Трунин. "--- Самара: CамГТУ"--~СамВен, 1997. "--- 308~с.


\item Справочник по плавкости систем из безводных неорганических солей [Текст]~/  Под ред. Н.\,К.~Воскресенской. "--- М.--Л.: АН СССР, 1961. "--- T.~2: Системы тройные и~более сложные. "--- С.~347--352.

\item Термические константы веществ [Текст]: Справочник в 10 выпусках / Под ред. акад. В.\,П.~Глушко. "---  М.: АН СССР, 1982.

\item Справочник по плавкости систем из безводных неорганических солей [Текст]~/ Под ред. Н.\,К.~Воскресенской. "--- М.--Л.: АН СССР, 1961. "--- T.~1: Двойные системы. "--- 845~с.


\item \textit{Зайдель, А.\,Н.} Ошибки измерения физических величин [Текст]~/  А.\,Н.~Зайдель. "---  Л.: Наука, 1974. "--- 108~с.



\item \textit{Посыпайко, В.\,И.} Применение ЭВМ при статистической обработке экспериментальных данных по нонвариантным точкам диаграмм состояния солевых систем [Текст]~/  
В.\,И.~Посыпайко, В.\,Н.~Первикова, В.\,Я.~Волков, Б.\,В.~Стратилатов~/\!/
Докл. АН СССР. "--- 1976. "--- Т.~223, №~3. "--- С.~669--671.



\item \textit{Трунин, А.\,С.} Визуально"=политермический метод [Текст]~/
А.\,С.~Трунин, Д.\,Г.~Петрова. "--- Куйбышев: КПтИ, 1977. "--- 93~с. "--- Деп. в ВИНИТИ 16.12.78 \No~548--78.


\item \textit{Егунов, В.\,П.} Введение в термический анализ [Текст]~/
В.\,П.~Егунов. "---  Самара, 1996. "--- 270~с.
\end{thebibliography}

\lastdate{15.02.2008}

\vspace{10mm}

\makeengtitle


%\end{document}
